\textbf{1. Kiến thức cần nhớ}
\begin{itemize}
\item Cho tam giác $ABC$ vuông tại $A$, xét góc nhọn $B = \alpha$:
\begin{itemize}
\item $\sin \alpha = \dfrac{\text{cạnh đối}}{\text{cạnh huyền}} = \dfrac{AC}{BC}$.
\item $\cos \alpha = \dfrac{\text{cạnh kề}}{\text{cạnh huyền}} = \dfrac{AB}{BC}$.
\item $\tan \alpha = \dfrac{\text{cạnh đối}}{\text{cạnh kề}} = \dfrac{AC}{AB}$.
\item $\cot \alpha = \dfrac{\text{cạnh kề}}{\text{cạnh đối}} = \dfrac{AB}{AC}$.
\end{itemize}
\item Mẹo nhớ: \textit{Sin đi học, Cứ khóc hoài, Thôi đừng khóc, Có kẹo đây}.
\item \textbf{Nhận xét}: Với góc nhọn $\alpha$ bất kì, ta luôn có:
$$\boxed{0 < \sin \alpha < 1;\quad 0 < \cos \alpha < 1;\quad \tan \alpha > 0;\quad \cot \alpha > 0}$$
\item \textbf{Tỉ số lượng giác của hai góc phụ nhau} ($\alpha + \beta = 90^\circ$):
$$\boxed{\sin \alpha = \cos \beta;\quad \cos \alpha = \sin \beta;\quad \tan \alpha = \cot \beta;\quad \cot \alpha = \tan \beta}$$
\item \textbf{Các hệ thức lượng giác cơ bản}:
$$\boxed{\tan \alpha = \dfrac{\sin \alpha}{\cos \alpha}};\quad \boxed{\cot \alpha = \dfrac{\cos \alpha}{\sin \alpha}};\quad \boxed{\tan \alpha \cdot \cot \alpha = 1}$$
$$\boxed{\sin^2 \alpha + \cos^2 \alpha = 1};\quad \boxed{1 + \tan^2 \alpha = \dfrac{1}{\cos^2 \alpha}};\quad \boxed{1 + \cot^2 \alpha = \dfrac{1}{\sin^2 \alpha}}$$
\end{itemize}

\textbf{2. Phương pháp giải các dạng}

%% DANG 9H4?11-1
\textbf{Dạng 1. Tính tỉ số lượng giác của góc nhọn}
\begin{enumerate}
\item Áp dụng định lí Pythagore để tìm độ dài cạnh còn thiếu của tam giác vuông: $BC^2 = AB^2 + AC^2$.
\item Viết công thức định nghĩa tỉ số lượng giác của góc nhọn cần tính.
\item Thay số và rút gọn phân số hoặc căn thức.
\end{enumerate}

%% DANG 9H4?11-2
\textbf{Dạng 2. Tỉ số lượng giác của hai góc phụ nhau}
\begin{enumerate}
\item Nhận diện hai góc có tổng số đo bằng $90^\circ$ (ví dụ: $20^\circ$ và $70^\circ$, $35^\circ$ và $55^\circ$).
\item Chuyển đổi tỉ số: $\sin 70^\circ = \cos 20^\circ$, $\tan 55^\circ = \cot 35^\circ$.
\item Thay thế để triệt tiêu hoặc thu gọn biểu thức.
\end{enumerate}

%% DANG 9H4?11-3
\textbf{Dạng 3. Tìm góc khi biết tỉ số lượng giác, so sánh}
\begin{enumerate}
\item Sử dụng máy tính cầm tay (nhấn phím $\text{SHIFT} \rightarrow \sin / \cos / \tan$) để tìm số đo góc xấp xỉ (đổi ra độ, phút).
\item So sánh hai góc nhọn: Khi góc $\alpha$ tăng từ $0^\circ$ đến $90^\circ$:
\begin{itemize}
\item $\sin \alpha$ và $\tan \alpha$ tăng.
\item $\cos \alpha$ và $\cot \alpha$ giảm.
\end{itemize}
\end{enumerate}

%% DANG 9H4?11-4
\textbf{Dạng 4. Rút gọn, chứng minh biểu thức lượng giác}
\begin{enumerate}
\item Sử dụng các hằng đẳng thức lượng giác: $\sin^2\alpha + \cos^2\alpha = 1$, $\tan\alpha \cdot \cot\alpha = 1$.
\item Đưa tất cả về $\sin$ và $\cos$ bằng cách thay $\tan\alpha = \dfrac{\sin\alpha}{\cos\alpha}$, $\cot\alpha = \dfrac{\cos\alpha}{\sin\alpha}$.
\item Quy đồng mẫu số và rút gọn.
\end{enumerate}
